Congressional apportionment is a mathematical problem with a constitutional
mandate and a century of contested solutions.
The choice of method is not arbitrary: it determines which states gain
or lose representation for a decade, and it determines whether the
resulting allocations are susceptible to counterintuitive paradoxes
that undermine public confidence in the process.

This paper has established the foundational framework for the J-track.
The divisor vs.\ quota taxonomy divides the five major methods into
two families with fundamentally different mathematical properties.
The three fairness criteria --- population pair test, quota rule,
and paradox immunity --- capture different values that cannot
simultaneously be satisfied: the Balinski-Young impossibility theorem
guarantees a tradeoff.

Congress's 1941 choice of Huntington-Hill reflects a specific resolution
of this tradeoff: paradox immunity and near-proportional representation
at the cost of occasional quota-rule violations.
The Supreme Court in \textit{Montana} confirmed that Congress has broad
discretion in choosing among apportionment methods and that the choice
need not be optimal under any single criterion.

The current \textsc{bisect-apportion} crate makes the divisor-method portion of
this framework computationally accessible: it implements Huntington-Hill,
Webster, Adams, and Jefferson with deterministic priority-queue logic.
Hamilton support, a standalone CLI, and checked-in Census-reference verification
fixtures remain future hardening work before the crate should be described as a
standalone official-output verifier.

\paragraph{Future directions.}
The J-track does not address the apportionment of state legislative
chambers, which may use different methods (many states use Webster
or Hamilton for their own chambers), or the apportionment of seats
in the Electoral College.
Extending \textsc{bisect-apportion} to these contexts, and to
non-U.S.\ parliamentary systems where d'Hondt and Sainte-Lagu\"e
are used for party-list PR, would expand the crate's utility
for comparative electoral systems research.
